\documentclass[11pt,a4paper]{article}

\usepackage{iftex}
\RequirePDFTeX
\usepackage{cmap}
\usepackage[T2A]{fontenc}
\usepackage[utf8]{inputenc}
\usepackage[english,russian]{babel}
\usepackage{PTSerif}
\usepackage{amsmath,amssymb,microtype}
\usepackage[a4paper,left=23mm,right=23mm,top=20mm,bottom=21mm]{geometry}
\usepackage{xcolor,enumitem,fancyhdr,hyperref}

\definecolor{accent}{HTML}{187AA5}
\definecolor{ink}{HTML}{302B28}
\definecolor{muted}{HTML}{6D6560}
\definecolor{tint}{HTML}{EDF4F6}
\definecolor{rulecolor}{HTML}{A9CADA}
\hypersetup{unicode=true,colorlinks=true,urlcolor=accent,linkcolor=accent,
  pdftitle={Домашнее задание 5. Околокруговое приближение. Алгоритмы перехода. Базис-вектор.},
  pdfauthor={Чиняев Владимир Николаевич}}
\urlstyle{same}
\color{ink}
\setlength{\parindent}{0pt}
\setlength{\parskip}{3pt}
\linespread{1.08}
\setlist[enumerate]{leftmargin=*,label=\textcolor{accent}{\arabic*.},
  itemsep=3pt,topsep=4pt,parsep=0pt}
\setlist[itemize]{leftmargin=*,itemsep=3pt,topsep=3pt,parsep=0pt}
\setlength{\emergencystretch}{1em}
\pagestyle{fancy}
\fancyhf{}
\renewcommand{\headrulewidth}{0pt}
\renewcommand{\footrulewidth}{0.4pt}
\renewcommand{\footrule}{\hbox to\headwidth{\color{rulecolor}\leaders\hrule height\footrulewidth\hfill}}
\fancyfoot[L]{\footnotesize\color{muted}Введение в маневрирование КА\quad /\quad ДЗ 5}
\fancyfoot[R]{\footnotesize\color{muted}\thepage}

\newcommand{\taskheading}[2]{%
  \par\vspace{5pt}{\large\bfseries\color{accent}Задание #1. #2}\par\smallskip}
\newcommand{\hint}[1]{%
  \par\nobreak\smallskip\noindent
  \begingroup
  \setlength{\fboxsep}{7pt}%
  \setlength{\fboxrule}{0.4pt}%
  \fcolorbox{rulecolor}{tint}{%
    \begin{minipage}{\dimexpr\linewidth-2\fboxsep-2\fboxrule\relax}
    \small\textbf{\color{accent}Подсказка.} #1\par
    \end{minipage}}%
  \endgroup\par\smallskip}

\begin{document}

{\small\color{accent}ВВЕДЕНИЕ В МАНЕВРИРОВАНИЕ КА}\par
\vspace{3pt}
{\huge Домашнее задание 5}\par
{\large Околокруговое приближение.\\Алгоритмы перехода. Базис-вектор.}\par
{\small\color{muted}К лекции 6}
\par{\small\bfseries\color{accent}Срок сдачи: 2 ноября 2026 года, 23:59}
\vspace{4pt}
{\color{rulecolor}\hrule}
\vspace{5pt}

Рассмотрите компланарный и пространственный переходы между близкими
орбитами в околокруговом приближении. Задания 1--3 обязательны,
задание 4* --- по желанию. Решение оформите в предоставленном Jupyter Notebook.

\taskheading{1}{Чтение методички}

В методичке \texttt{near-circular-transfer.pdf} прочитайте разделы
3.1, 3.2.2, 3.3.1 и 3.4: постановку задачи перехода, алгоритм для
непересекающихся компланарных орбит, узловое решение и порядок выбора алгоритмов.
Повторите достаточные условия оптимальности с базис-вектором.
Остальные пространственные случаи просмотрите на уровне геометрии
годографа, областей применимости и итоговых алгоритмов.
Подробный вывод невырожденного и особого решений повторять не требуется.

Дайте краткие ответы без подробного вывода формул:
\begin{enumerate}
\item Почему для непересекающихся компланарных орбит место первого импульса
можно выбирать свободно? Как это связано с годографом базис-вектора?
\item В методичке приведены два условия применимости узлового решения.
Какое из них обеспечивает неотрицательность модулей импульсов,
а какое --- ограничение модуля базис-вектора на всём обороте?
\end{enumerate}

\begingroup
\setlength{\fboxsep}{9pt}
\colorbox{tint}{\begin{minipage}{\dimexpr\linewidth-2\fboxsep\relax}
\textbf{Общие исходные данные}
\[
 \mu=398600{,}44158\ \text{км}^3/\text{с}^2,\qquad
 r_0=7000\ \text{км}.
\]
\end{minipage}}
\endgroup

Используйте обозначения и нормировку из методички.
Начало маневрирования, длительность перехода и конечное положение
аппарата на орбите свободны.

\vspace{4pt}
\textbf{Работа с ноутбуком и самопроверка}

В \texttt{hw\_5.ipynb} готовы константы, таблицы, графики и чекеры
\texttt{check\_task2}, \texttt{check\_task3}, \texttt{check\_star}.
Заполните отмеченные переменные и текстовые ячейки своими решениями.
Использование чекеров не требуется.

Передавайте углы в радианах, компоненты --- в нормированном виде.
Каждая строка массива импульсов содержит
$(\Delta\widehat V_r,\Delta\widehat V_t,\Delta\widehat V_z)$.
Множители передаются в порядке $(\lambda_1,\lambda_2,\lambda_3,\lambda_4,\lambda_5)$.
Расчёты выполняйте без промежуточного округления.
Чекеры проверяют численные ответы; рассуждения и доказательства
оформите самостоятельно.

\newpage
{\small\color{muted}Домашнее задание 5\hfill К лекции 6}\par
\vspace{5pt}

\taskheading{2}{Разные манёвры с одинаковым расходом}

Заданы изменения элементов
\[
 \Delta a=0{,}004,\qquad
 \Delta e_x=0{,}001,\qquad
 \Delta e_y=0{,}001,\qquad \Delta i=0.
\]
Ось отсчёта углов фиксирована в общей плоскости орбит;
компоненты векторов эксцентриситета обеих орбит заданы относительно этой оси.
\begin{enumerate}
\item Определите взаимное расположение орбит и выберите алгоритм
перехода. Объясните, возможен ли переход одним импульсом.
\item Рассчитайте три двухимпульсных перехода, задавая положение
первого импульса:
\[
 \varphi_1=0^\circ,\quad120^\circ,\quad240^\circ.
\]
Для каждого найдите $\varphi_2$, обе трансверсальные компоненты
и суммарную характеристическую скорость.
Радиальные и нормальные компоненты положите равными нулю.
Представьте таблицу: углы в градусах, компоненты и суммарная скорость в м/с.
Углы сначала найдите по модулю $2\pi$, затем запишите с учётом
последовательности импульсов. При необходимости добавьте $2\pi$ к $\varphi_2$.
\item Для каждого перехода подстановкой проверьте три граничных условия:
\[
\begin{aligned}
 2\sum_j\Delta\widehat V_{tj}&=\Delta a,\\
 \sum_j2\Delta\widehat V_{tj}\cos\varphi_j&=\Delta e_x,\\
 \sum_j2\Delta\widehat V_{tj}\sin\varphi_j&=\Delta e_y.
\end{aligned}
\]
\item Найдите \emph{один общий базис-вектор}, подтверждающий оптимальность
всех трёх решений. Аналитически проверьте достаточные условия
для каждого перехода.
\item Объясните, почему места приложения и величины отдельных импульсов
различаются, а суммарная характеристическая скорость совпадает.
Что геометрия годографа говорит о возможности выбора места первого импульса?
\end{enumerate}

В заданиях 2 и 3 используйте базис-вектор задачи перехода:
\[
\begin{aligned}
 s_r&=\lambda_3\sin\varphi-\lambda_2\cos\varphi,\\
 s_t&=2\lambda_1+2\lambda_3\cos\varphi+2\lambda_2\sin\varphi,\\
 s_z&=\lambda_4\sin\varphi+\lambda_5\cos\varphi.
\end{aligned}
\]
Здесь $\lambda_6=0$. В компланарном переходе положите
$\lambda_4=\lambda_5=0$.

\hint{Используйте алгоритм для непересекающихся орбит со свободным
углом первого импульса. Второй угол восстанавливайте по паре синуса
и косинуса с помощью \texttt{atan2}; один арккосинус не определяет четверть.
При поиске общего базис-вектора сопоставьте знаки трансверсальных
компонент всех импульсов. Компоненты базис-вектора записаны
в местных орбитальных осях.}

\newpage
{\small\color{muted}Домашнее задание 5\hfill К лекции 6}\par
\vspace{5pt}

\taskheading{3}{Узловой переход и его оптимальность}

Заданы
\[
 \Delta a=0{,}001,\qquad
 \Delta e_x=0{,}0024,\qquad
 \Delta e_y=0{,}0006,\qquad
 \Delta i=0{,}0018\ \text{рад}.
\]
Ось $X$ направлена к восходящему узлу конечной плоскости относительно
начальной; $\Delta e_x,\Delta e_y$ заданы в этих узловых осях.
\begin{enumerate}
\item Проверьте применимость узлового алгоритма.
\item Найдите положения импульсов, их модули и все три компоненты
каждого импульса. Рассчитайте суммарную характеристическую скорость.
Представьте таблицу с углами в градусах и скоростями в м/с.
\item Подстановкой проверьте все пять граничных условий:
\[
\begin{aligned}
 2\sum_j\Delta\widehat V_{tj}&=\Delta a,\\
 \sum_j(\Delta\widehat V_{rj}\sin\varphi_j
       +2\Delta\widehat V_{tj}\cos\varphi_j)&=\Delta e_x,\\
 \sum_j(-\Delta\widehat V_{rj}\cos\varphi_j
       +2\Delta\widehat V_{tj}\sin\varphi_j)&=\Delta e_y,\\
 \sum_j\Delta\widehat V_{zj}\sin\varphi_j&=0,\\
 \sum_j\Delta\widehat V_{zj}\cos\varphi_j&=\Delta i.
\end{aligned}
\]
\item Определите множители $\lambda_1,\ldots,\lambda_5$ и запишите
соответствующий базис-вектор.
Поясните, какие условия использованы для нахождения множителей.
\item Постройте пространственный годограф с единичной сферой
и график $\|\boldsymbol s(\varphi)\|$ при $0\leq\varphi\leq2\pi$.
Отметьте положения импульсов на обоих графиках.
\item Аналитически докажите $\|\boldsymbol s(\varphi)\|\leq1$
при всех $\varphi$. Проверьте единичный модуль в точках импульсов
и сонаправленность базис-вектора каждому импульсу.
По достаточным условиям обоснуйте глобальную оптимальность
среди манёвров с произвольным числом импульсов в околокруговом приближении.
\item Объясните, почему противоположные направления импульсов
в местных орбитальных осях допускают разные модули.
Как изменение $\Delta a$ внутри области применимости влияет
на распределение импульсов, их направления и суммарную величину?
\end{enumerate}

\hint{В узлах синусы и косинусы принимают простые значения.
Сначала найдите сумму модулей и общее единичное направление,
затем распределите модули по условию на $\Delta a$.
Для проверки базис-вектора запишите его как сумму постоянных
векторов при $\cos\varphi$ и $\sin\varphi$ и исследуйте квадрат модуля.
График иллюстрирует результат; доказательство должно охватывать
все углы между узлами сетки.}

\newpage
{\small\color{muted}Домашнее задание 5\hfill К лекции 6}\par
\vspace{5pt}

\taskheading{4*}{Граница узлового и особого случаев}
\textit{Необязательное задание.}

Сохраните $\Delta a,\Delta e_x,\Delta e_y$ из задания 3,
а взаимный угол плоскостей изменяйте в диапазоне
\[
 0{,}0003\leq\Delta i\leq0{,}0018\ \text{рад}.
\]
\begin{enumerate}
\item По условиям классификации определите области применимости
узлового и особого алгоритмов на заданном отрезке.
Найдите границу между ними. Для особого случая проверьте оба ограничения.
\item Постройте на одном графике характеристические скорости
по формулам узлового и особого случаев.
Выделите участки, на которых каждая формула применима,
и постройте итоговую зависимость оптимальной стоимости перехода.
По горизонтальной оси отложите $\Delta i$ в радианах,
по вертикальной --- $\Delta V_\Sigma$ в м/с.
\item Аналитически проверьте совпадение стоимостей на границе.
\item Исследуйте годограф \emph{узлового} базис-вектора
при $\Delta i=0{,}0003$~рад, на найденной границе
и при $\Delta i=0{,}0018$~рад. Для каждого значения постройте
годограф с единичной сферой и график модуля.
Объясните изменение геометрии и проверьте условие
$\|\boldsymbol s(\varphi)\|\leq1$.
\item Можно ли выбирать алгоритм, подставляя данные в обе формулы
стоимости и принимая меньший результат?
Обоснуйте ответ сравнением формул и условиями их применимости.
\end{enumerate}

\hint{Различайте вычисление выражения для стоимости и существование
оптимального манёвра данного типа.
На общей границе могут совпадать и стоимость, и геометрия годографа.
Для сравнения формул удобно рассмотреть разность квадратов
нормированных стоимостей. Полный расчёт трёхимпульсного особого перехода
в этом задании не требуется.}

\vspace{5pt}
{\color{rulecolor}\hrule}
\vspace{5pt}
\textbf{Что представить}

Сдайте один ноутбук с полными решениями заданий 2 и 3
и, по желанию, задания 4*: формулами, таблицами импульсов,
проверкой граничных условий, годографами и графиками модуля
базис-вектора, обоснованием оптимальности и ответами на вопросы.
В задании 1 приведите краткие ответы на оба вопроса.
Использование чекеров и сохранение результатов их вызова не требуются.

Формулы и рассуждения оформите в текстовых ячейках, расчёты и рисунки ---
в кодовых. Перед отправкой перезапустите ядро, выполните ячейки решения по порядку
и сохраните ноутбук с результатами.

Название файла: \texttt{hw\_5\_name\_surname.ipynb}.
Вместо \texttt{name\_surname} укажите свои имя и фамилию латиницей.

\end{document}
